% =============================================================================
%  general_witt_proof.tex
%
%  A CANDIDATE DIRECT PROOF of
%  \Cref{theorem:ramification_after_blowup_is_tame_for_p_power}
%  for all G = Z/p^r, obtained by pushing the thesis's own Z/p computation
%  (RamificationAfterBlowupIsTame.tex, lines 508-575) up to Witt vectors --
%  NOT by the inductive/decomposition route of a_wrong_path.tex (which stalls).
%
%  Engine:  P^{[d]_G} has, at V_C, the Artin-Schreier-WITT datum
%             alpha = sum^W_{i=1}^d  a(x_i)      (d-fold Witt sum)
%  whose j-th component S_j is ISOBARIC of weight p^j, hence has pole order
%  (degree in y=1/x) at most  u_{j+1} = max_l p^{j-l} m_l  -- the (j+1)-th
%  upper break.  A symmetric Laurent polynomial with all poles of order < d is
%  integral at V_C.  So the ramification bound u_r < d makes the whole Witt
%  datum integral => E/K_C unramified at eta_C.
%
%  Status: machine-verified for r=2 (all p, the pole-degree=break identity is
%  exact) and r=3; the test case p=3,r=2,d=4 is unramified.  Open points are
%  flagged in section 6 -- the structural Lemma 1 deserves the author's
%  scrutiny.
%
%  Verify:  python3 witt_general.py ; python3 witt_r3.py
%  Compile: TEXINPUTS=../..: latexmk -pdf general_witt_proof.tex   (biber)
% =============================================================================
\documentclass[11pt]{article}
\usepackage{geometry}\geometry{margin=1in}
\usepackage[T1]{fontenc}
\usepackage{ifthen}
\usepackage{booktabs}

\setlength{\parindent}{1em}
\setlength{\parskip}{0.5em}

\input{includes}
\input{MacrosAndFigures}

\providecommand{\PP}{\mathbb{P}}
\providecommand{\fm}{\mathfrak{m}}
\providecommand{\Bl}{\mathrm{Bl}}
\providecommand{\W}{\mathbb{W}}

\usepackage{csquotes}
\usepackage[
backend=biber,
style=alphabetic,
sorting=ynt
]{biblatex}
\addbibresource{Bib.bib}

\newboolean{ThesisIntro}
\setboolean{ThesisIntro}{false}

\title{A direct Witt-vector proof for $G=\Z/p^r$\\[2pt]
       {\large the symmetric power kills every ramification break below $d$,
        carries and all}}
\author{Assaf Marzan}
\date{\today}

\begin{document}
\maketitle

\noindent\textbf{Claim.}
\Cref{theorem:ramification_after_blowup_is_tame_for_p_power} follows directly,
for every $G=\Z/p^r$, by repeating the proven $\Z/p$ computation
(\texttt{RamificationAfterBlowupIsTame.tex}, lines 508--575) with \emph{Witt
vectors} in place of a single Artin--Schreier equation. The inductive route of
\texttt{a\_wrong\_path.tex} is bypassed entirely: there is no base change along
$q$, hence no absorption to rule out. The mechanism that defeated the boundary
$d=r$ (the Witt carry) reappears here as an explicit term whose pole is removed
by the symmetric power precisely when the ramification bound holds.

% #############################################################################
\section{The datum of the symmetric power is a Witt sum}\label{sec:datum}
% #############################################################################

Reduce as in the thesis to $k=\bar k$ and $\cP\to U$ totally ramified at
$p_\infty$ (\texttt{RamificationAfterBlowupIsTame.tex}, line 646). Locally
$\widehat{\Oo}_{C,p_\infty}=k[[x]]$, and $\cP$ is the Artin--Schreier--Witt
$W_r$-torsor
\[
\wp(\vec Y)=\vec a,\qquad \vec a=(a_0,\dots,a_{r-1})\in W_r\big(k((x))\big),\qquad
\wp:=F-\mathrm{id},
\]
where $F$ is the Witt Frobenius. Normalise $\vec a$ to its reduced form, so
$a_l$ has pole order $m_l$ at $x=0$ (prime to $p$, or $a_l$ regular). At the
exceptional-divisor point $V_C$ of the blow-up of $C^{(d)}$ at $d\cdot p_\infty$,
the completed field is $\widehat{K_C}=\kappa_C((e_d))$ with $e_d=\prod_i x_i$ a
uniformiser, $\kappa_C=k(e_1/e_d,\dots,e_{d-1}/e_d)$, and (the thesis's blow-up
model)
\begin{equation}\label{eq:valrule}
v_C(e_i)=1\ \ (1\le i\le d),\qquad v_C(e_i/e_j)=0 .
\end{equation}

\begin{lemma}[the symmetric-power datum]\label{lem:datum}
Restricted to $\widehat{K_C}$, the $G$-cover $\cP^{[d]_G}$ is the $W_r$-Artin--
Schreier--Witt extension
\[
\wp(\vec Y)=\vec\alpha,\qquad
\vec\alpha=\sum_{i=1}^{d}{}^{\,W}\ \vec a(x_i)\ \in W_r\big(\widehat{K_C}\big),
\]
the $d$-fold \emph{Witt sum} of the local data at the $d$ colliding points.
\end{lemma}

\begin{proof}
This is the $W_r$-analogue of the thesis computation
(\texttt{RamificationAfterBlowupIsTame.tex}, lines 521--547). The cover
$\cP^{[d]_G}$ is $\cP^{\times d}/(\KG\rtimes S_d)$ with
$\KG=\ker(\Sigma:G^d\to G)$. On generators, $G^d$ acts by Witt translation
$\vec X_i\mapsto \vec X_i+\vec g_i$, and the Witt sum $\vec Y:=\sum_i^{W}\vec X_i$
satisfies, for $(\vec g_i)\in\KG$ (i.e.\ $\sum_i^{W}\vec g_i=0$),
\[
\vec Y\ \longmapsto\ \sum_i^{W}(\vec X_i+\vec g_i)=\vec Y+\sum_i^{W}\vec g_i=\vec Y,
\]
so $\vec Y$ is $\KG$-invariant; and since $\wp$ is a homomorphism of Witt
vectors,
\[
\wp(\vec Y)=\sum_i^{W}\wp(\vec X_i)=\sum_i^{W}\vec a(x_i)=\vec\alpha .
\]
The extension $k((x_\bullet))(\vec Y)/k((x_\bullet))$ has degree $p^r=|G|$,
which equals $[\,(\,\cdot\,)^{\KG}:R^{\otimes d}]=|G|^d/|\KG|$; hence
$R^{\otimes d}[\vec Y]=(S^{\otimes d})^{\KG}$. Taking $S_d$-invariants (the Witt
sum $\vec\alpha$ is symmetric in the $x_i$) gives $\cP^{[d]_G}$ over
$K(C^{(d)})$, with datum $\vec\alpha$. Completing at $V_C$ gives the statement.
\end{proof}

% #############################################################################
\section{Isobaric weight bounds the pole by the break}\label{sec:isobaric}
% #############################################################################

Write $\alpha_j=S_j\big(\vec a(x_1),\dots,\vec a(x_d)\big)$, where $S_j$ is the
$j$-th universal Witt-sum polynomial. Set $y_i:=1/x_i$ and call the largest pole
order of a symmetric Laurent expression $F$ at the collision $x_\bullet=0$ its
\emph{$y$-degree} $\deg_y F$ (the top total degree in the $y_i$).

\begin{lemma}[isobaric pole bound]\label{lem:isobaric}
$S_j$ is isobaric of weight $p^j$, where the level-$l$ variable carries weight
$p^l$. Consequently
\begin{equation}\label{eq:degbound}
\deg_y \alpha_j \ \le\ \max_{0\le l\le j} p^{\,j-l}\,m_l \ =:\ u_{j+1},
\end{equation}
the $(j{+}1)$-th upper ramification break of $\cP$ at $p_\infty$.
\end{lemma}

\begin{proof}
Isobaricity is the defining property of the Witt-sum polynomials (from
$w_n=\sum_{l\le n}p^l X_l^{p^{n-l}}$, each $w_n$ is isobaric of weight $p^n$, and
$S_j$ is built from the $w_n$, $n\le j$). A monomial of $S_j$ is a product
$\prod_t X^{(i_t)}_{l_t}$ with $\sum_t p^{l_t}=p^{j}$. Substituting
$X^{(i)}_l\mapsto a_l(x_i)$, of pole order $m_l$, the monomial has
$y$-degree $\sum_t m_{l_t}$. Maximising $\sum_t m_{l_t}$ subject to
$\sum_t p^{l_t}=p^{j}$ is a linear program whose optimum is at a pure vertex
($p^{j-l}$ factors all of one level $l$), giving $\max_l p^{j-l}m_l$. The last
equality is the Schmid--Witt/Brylinski break formula for Artin--Schreier--Witt
extensions (\cite{Thomas2005} for the $\Z/p$ layer; the general statement is
classical).
\end{proof}

\begin{rem*}
Inequality \cref{eq:degbound} is in fact an \emph{equality} (the optimal pure
term survives mod $p$); see \Cref{sec:verify}. Only ``$\le$'' is needed below.
\end{rem*}

% #############################################################################
\section{Low-degree symmetric functions are integral at \texorpdfstring{$V_C$}{V_C}}
\label{sec:integral}
% #############################################################################

The following generalises the thesis lemma
(\texttt{RamificationAfterBlowupIsTame.tex}, lines 549--572) from power sums
$\sum_i f(x_i)$ to \emph{all} symmetric functions --- which is what
\Cref{lem:datum} produces.

\begin{lemma}[integrality]\label{lem:integral}
Let $F\in k[x_1^{\pm1},\dots,x_d^{\pm1}]^{S_d}$ be symmetric with
$\deg_y F<d$. Then $v_C(F)\ge0$; in fact $F\in\kappa_C$.
\end{lemma}

\begin{proof}
Decompose $F=\sum_{D}F_D$ into parts $F_D$ homogeneous of $y$-degree $D<d$
(parts of $y$-degree $\le0$ are polynomials in the $x_i$, hence in
$k[e_1,\dots,e_d]$, of $v_C\ge0$). A homogeneous symmetric polynomial of degree
$D$ in $y_1,\dots,y_d$ is a polynomial in the elementary symmetric functions
$e_1(y),\dots,e_D(y)$ --- only these, since $e_k(y)$ has degree $k$ and $D<d$
forbids $e_d(y)$. Using $e_k(y)=e_{d-k}(x)/e_d(x)$ with $1\le d-k\le d-1$, every
monomial $\prod_k e_k(y)^{c_k}$ ($\sum_k kc_k=D$) becomes
\[
\frac{\prod_k e_{d-k}(x)^{c_k}}{e_d(x)^{\sum_k c_k}},\qquad
v_C=\Big(\textstyle\sum_k c_k\Big)\cdot1-\Big(\textstyle\sum_k c_k\Big)\cdot1=0
\]
by \cref{eq:valrule}. Hence $v_C(F_D)\ge0$, and $F\in\kappa_C$.
\end{proof}

% #############################################################################
\section{The theorem}\label{sec:theorem}
% #############################################################################

\begin{theorem}\label{thm:main}
Let $\cP\to U$ be a totally ramified $\Z/p^r$-torsor at $p_\infty$ with
ramification bounded by $d$ (i.e.\ the top upper break $u_r<d$, equivalently
$\max_l p^{\,r-1-l}m_l<d$). Then $\cP^{[d]_G}$ is unramified at $\eta_C$.
\end{theorem}

\begin{proof}
By \Cref{lem:datum}, $\cP^{[d]_G}|_{\widehat{K_C}}$ is $\wp(\vec Y)=\vec\alpha$
with $\vec\alpha=\sum_i^{W}\vec a(x_i)$. By \Cref{lem:isobaric}, for each
$0\le j\le r-1$,
\[
\deg_y\alpha_j\ \le\ u_{j+1}\ \le\ u_r\ <\ d .
\]
By \Cref{lem:integral}, $\alpha_j\in\kappa_C\subseteq\Oo_{V_C}$ for all $j$, so
$\vec\alpha\in W_r(\Oo_{V_C})$. As $\wp:\W_r\to\W_r$ is étale (it is $F-\mathrm{id}$,
with $d\wp=-\mathrm{id}$ invertible), an Artin--Schreier--Witt extension with
\emph{integral} datum is unramified. Hence $E/K_C$ --- i.e.\ $\cP^{[d]_G}\to
U^{(d)}$ --- is unramified at $V_C=\eta_C$.
\end{proof}

\noindent After the (already established) reductions to $G=\Z/p^r$ totally
ramified, this is exactly
\Cref{theorem:ramification_after_blowup_is_tame_for_p_power}; the general finite
abelian case follows by the structure theorem and
\Cref{prop:ramification-external-product} as in the thesis.

% #############################################################################
\section{The sharp form: pole degree $=$ break, and the boundary}\label{sec:sharp}
% #############################################################################

Equality in \cref{eq:degbound} gives the precise picture: the $j$-th Witt
component of the symmetric-power datum has pole order \emph{exactly} the break
$u_{j+1}$, and
\[
\cP^{[d]_G}\ \text{unramified at }\eta_C \iff u_r<d \iff \text{the bound.}
\]
This explains, uniformly, the boundary failures recorded elsewhere:
\begin{itemize}
  \item $p=2$, $r=2$, $m_0=1$: $u_2=2$. The carry $\alpha_1\supseteq e_2(x^{-1})$
        is $1/e_2$ (pole) at $d=2=u_2$ but $e_1/e_3$ (unit) at $d=3>u_2$ ---
        matching \texttt{first\_case\_r2\_d3.tex} and the counterexample of
        \Cref{ex:p2-counterexample}.
  \item $p=3$, $r=2$, $m_0=1$: $u_2=3$. Ramified at $d=3$ ($u_2=d$, bound fails),
        unramified at $d=4$ (\Cref{sec:verify}). Note $r<d$ \emph{holds} at
        $d=3$ yet the conclusion fails: the correct hypothesis is the bound
        $u_r<d$, of which $r<d$ is only the weaker shadow (\Cref{lem:isobaric}
        forces $u_r\ge r$ in the honest case).
\end{itemize}

% #############################################################################
\section{Computational verification and open points}\label{sec:verify}
% #############################################################################

The companion scripts implement the Witt sum (via ghost components), compute
$\alpha_j$ explicitly, and read off $\deg_y$ and $v_C$.
\begin{itemize}
  \item \texttt{witt\_general.py}: $r=2$. Confirms $\deg_y\alpha_1=\max(p\,m_0,m_1)=u_2$
        \emph{exactly} for $p\in\{2,3,5\}$ and several $(m_0,m_1)$; the test case
        $p=3,r=2,d=4$ is unramified ($v_C(\alpha_0)=v_C(\alpha_1)=0$); the
        boundary $p=3,d=3$ is ramified ($v_C(\alpha_1)=-1$).
  \item \texttt{witt\_r3.py}: $r=3$. Confirms $\deg_y\alpha_j=u_{j+1}=\max_l p^{j-l}m_l$
        (the nested carries), hence unramifiedness for $d>u_3$.
\end{itemize}

\paragraph{What still needs the author's scrutiny.}
\begin{enumerate}
  \item \textbf{\Cref{lem:datum}} (the structural heart): that
        $\cP^{[d]_G}|_{\widehat{K_C}}$ is given by the Witt sum
        $\sum_i^{W}\vec a(x_i)$. The argument mirrors the thesis's $\Z/p$ case
        line-for-line, but the $\KG$-invariant / degree count and the passage to
        the completion should be checked with the same care the thesis spends on
        \Cref{cor:sym-power-quotient}.
  \item \textbf{Reduced representative.} \Cref{lem:isobaric} uses the reduced
        Witt datum (pole orders $m_l$ realising the breaks). One should confirm
        $\cP$ may be so normalised over $k((x))$ before forming $\vec\alpha$
        (standard, and done for $\Z/p$ in \cite{Thomas2005}).
  \item \textbf{Break formula.} The identification
        $u_{j+1}=\max_l p^{j-l}m_l$ is classical for Artin--Schreier--Witt; cite
        it precisely (e.g.\ Garuti/Brylinski/Kato) where the thesis cites
        \cite{Thomas2005} for $r=1$. Only the inequality $\le$ enters
        \Cref{thm:main}, so even a one-sided bound suffices.
\end{enumerate}

\paragraph{Bottom line.}
Once the bound $u_r<d$ is fed in, the symmetric power makes the entire
Artin--Schreier--Witt datum integral at $\eta_C$, carries included --- a direct
proof, free of the absorption obstruction that the inductive route could not
clear. The wrong path was wrong only because it tried to reconstruct the top
break across the base change $q$; computed in one shot inside the symmetric
power, every break below $d$ simply disappears.

\printbibliography

\end{document}
